%======================= PROBLÉMATIQUE ET OBJECTIFS =========================
\chapter{Problématique et objectifs}

\section{Problématique}

La conduite d'une campagne \emph{drive-to-store} met en jeu des étapes
hétérogènes, souvent gérées par des outils distincts : définition de la
campagne et de son budget, collecte et validation de la base de magasins,
définition des zones de diffusion autour des points de vente, production de
créatives personnalisées, puis analyse des performances. Cette dispersion
soulève plusieurs difficultés :

\begin{itemize}
  \item \textbf{manque de centralisation} : les données (annonceurs,
        magasins, campagnes) sont ressaisies ou dupliquées d'un outil à
        l'autre ;
  \item \textbf{cohérence des données} : sans validation à l'import, une base
        de magasins peut contenir des erreurs (coordonnées manquantes, format
        d'URL invalide, doublons) qui faussent le ciblage ;
  \item \textbf{personnalisation à l'échelle} : produire une créative par
        magasin manuellement n'est pas viable dès que le réseau grandit ;
  \item \textbf{lisibilité du reporting} : les indicateurs doivent être
        filtrables et exportables pour être réellement exploitables ;
  \item \textbf{gestion des accès} : différents profils (administration, achat
        média, lecture seule) ne doivent pas avoir les mêmes droits.
\end{itemize}

La problématique centrale est donc la suivante :
\begin{quote}
\itshape
Comment concevoir une plateforme web unique qui centralise la création de
campagnes drive-to-store, le ciblage géographique des magasins, la génération
de créatives personnalisées et le reporting, tout en offrant une base
technique propre et évolutive ?
\end{quote}

\section{Objectifs du projet}

Pour répondre à cette problématique, le projet s'est fixé les objectifs
suivants :

\subsection{Objectifs fonctionnels}
\begin{enumerate}[label=\textbf{O\arabic*.}]
  \item Offrir une \textbf{authentification} et une \textbf{navigation
        protégée} par rôle (administration, achat média, lecture seule).
  \item Proposer un \textbf{tableau de bord} synthétisant les indicateurs
        clés et les campagnes récentes.
  \item Permettre la \textbf{gestion des campagnes} et leur \textbf{création
        guidée} via un assistant multi-étapes.
  \item Prendre en charge l'\textbf{import et la validation} d'une base de
        magasins (fichiers \texttt{.xlsx} / \texttt{.csv}).
  \item Assurer le \textbf{ciblage géographique} (\emph{geofencing}) sur une
        carte interactive.
  \item Générer des \textbf{créatives dynamiques (DCO)} et des \emph{landing
        pages} personnalisées par magasin.
  \item Fournir un \textbf{reporting} filtrable avec indicateurs, graphiques,
        carte, tableau et export CSV.
  \item Mettre à disposition un module d'\textbf{administration du compte}
        (annonceurs, utilisateurs, paramètres).
\end{enumerate}

\subsection{Objectifs techniques}
\begin{enumerate}[label=\textbf{T\arabic*.}]
  \item Construire une \textbf{architecture fullstack claire} (frontend
        React, backend FastAPI) communiquant via une API REST.
  \item Documenter l'API de manière \textbf{interactive} (Swagger/OpenAPI).
  \item Préparer une \textbf{base de données PostgreSQL} (modèles et
        migrations) sans bloquer le fonctionnement du prototype.
  \item Adopter un \textbf{flux de versionnement Git/GitHub} rigoureux
        (branches, \emph{pull requests}).
  \item Maintenir une \textbf{documentation} technique et fonctionnelle à jour.
\end{enumerate}

\section{Périmètre et hors-périmètre}

Le projet se concentre sur l'\textbf{expérience de gestion de campagne} et son
socle technique. Sont explicitement \textbf{hors périmètre} à ce stade : la
connexion à un réseau publicitaire réel (enchères RTB / OpenRTB), le
déploiement en production, l'intégration d'un moteur de DCO réel, et la mise en
place d'une authentification de niveau production. Ces éléments sont abordés
dans le chapitre \og Limites et perspectives \fg{}.
